\lysh{15275_SOLUTION}{1}
\textbf{ΛΥΣΗ}

\textbf{α) i.}
Αφού η ευθεία διέρχεται από το $M(2,1)$ και έχει συντελεστή διεύθυνσης $\lambda$, η εξίσωσή της είναι $y-1=\lambda(x-2)$ που γράφεται $y=\lambda x-2\lambda+1$.

\textbf{ii.}
Σε κάθε περίπτωση που ισχύει $\lambda \neq 0$ η ευθεία τέμνει και τους δύο άξονες και όταν δεν διέρχεται από την αρχή $O$, δηλαδή όταν $\lambda \neq \frac{1}{2}$, σχηματίζει τρίγωνο.

Επομένως, η ευθεία σχηματίζει τρίγωνο με τους άξονες μόνο όταν $\lambda \in \mathbb{R}-\left\{0, \frac{1}{2}\right\}$.

\textbf{β) i.}
Με $y=0$ στην εξίσωση της ευθείας, παίρνουμε $x=\frac{2\lambda-1}{\lambda}$, οπότε $A\left(\frac{2\lambda-1}{\lambda}, 0\right)$, ενώ με $x=0$ παίρνουμε $y=-2\lambda+1$, οπότε $B(0, -2\lambda+1)$. Επομένως τα μήκη των τμημάτων OA, OB είναι:
\[ (\mathrm{OA})=\left|\frac{2\lambda-1}{\lambda}\right| \quad \text{και} \quad (\mathrm{OB})=|-2\lambda+1|=|2\lambda-1| \]

\textbf{ii.}
Η ευθεία σχηματίζει με τους άξονες ισοσκελές τρίγωνο, μόνο όταν $(\mathrm{OA})=(\mathrm{OB})$. Είναι:
\[ (\mathrm{OA})=(\mathrm{OB}) \Leftrightarrow \left|\frac{2\lambda-1}{\lambda}\right|=|2\lambda-1| \Leftrightarrow |2\lambda-1|(|\lambda|-1)=0 \Leftrightarrow |\lambda|-1=0 \]
αφού $\lambda \neq \frac{1}{2}$. Άρα η ευθεία σχηματίζει ισοσκελές τρίγωνο, όταν $\lambda=-1$ ή $\lambda=1$.

\textbf{iii.}
Αν $\lambda=-1$, τότε $(\mathrm{OA})=3$ και $(\mathrm{OB})=3$, οπότε το εμβαδόν $(\mathrm{OAB})$ του τριγώνου $\mathrm{OAB}$ είναι
\[ (\mathrm{OAB})=\frac{9}{2} \]
Αν $\lambda=1$, τότε $(\mathrm{OA})=1$ και $(\mathrm{OB})=1$, οπότε το εμβαδόν του τριγώνου $\mathrm{OAB}$ είναι $(\mathrm{OAB})=\frac{1}{2}$

\textbf{Σχόλιο}

Το τρίγωνο $\mathrm{OAB}$ είναι ορθογώνιο και ισοσκελές, οπότε η γωνία που σχηματίζει η ευθεία (υποτείνουσα) με τον άξονα $x'x$ είναι $45^\circ$ ή $135^\circ$. Έτσι, έχουμε $\lambda=\varepsilon\phi 45^\circ=1$ ή $\lambda=\varepsilon\phi 135^\circ=-1$ που είναι οι τιμές που βρήκαμε παραπάνω.
\end{lysh}